% Discussion. Numbers: ../FINDINGS.md.
\section{Discussion}\label{sec:discussion}

\paragraph{The signal is real and not just the journal.} Citing a paper that has already been retracted is followed by a
markedly higher chance of retraction in the next three years, and the association survives comparison within journal
and publication year (odds ratio about 3). This turns the retrospective association reported by
\citet{lou2026beyond}, who compared citing articles with the global retraction rate, into a prospective one. Part of
the crude association does reflect journals: the crude risk ratio of 8.8 shrinks to an odds ratio of 3.0 within journal
and year (for an outcome this rare the two measures are nearly equal). But a threefold excess remains, so the signal is not explained by retraction waves within journals, and holding
reference-list length fixed does not change it. Their rates are also higher than ours (9.58\% at five retracted
references against 2.27\% here for non-review papers): they count retractions over the whole history of each citing
article and include articles retracted together with the papers they cite, whereas here only retractions within 2.75
years after $T_0$ count and every paper already retracted by $T_0$ is excluded.
The rate rises with the number of retracted references, and among non-review papers it keeps rising up to five or
more; citing a paper-mill retraction is the strongest single marker. Being cited by a retracted paper is at least as informative as citing one, which fits the
picture of interlinked clusters of problematic papers \citep{candal2024references}.

\paragraph{But it is a weak screen on its own.} A relative risk of nine on a base rate of 0.04\% still means that almost
all flagged papers are not retracted: about 1 in 110 of the top 1,000 by $S_1$. And 84\% of later retractions cite no
retracted paper at all, so the signal misses most cases. A ``Feet of Clay'' flag is therefore a reason to look, not
evidence of a problem. Many papers cite retracted work for good reasons, such as reviews, critiques, or citations made
before the retraction. Reviews are a clear case: they make up most of the papers with many retracted references and are almost
never retracted, so a screen should set them aside, which raised precision among the top 1,000 from 0.9\% to 2.2\%; among other papers the risk keeps rising with the number of
retracted references.

\paragraph{Timing matters.} The best screen depends on where the time point falls relative to retraction waves. Inside
the late-2023 wave, the journal's own retraction rate identified future retractions with a lift of about 90, and the
neighbourhood signal added little for paper-mill cases. Once the wave had passed, the journal rate carried almost no
information about who would be retracted next, while the neighbourhood signal kept its strength. A practical screen
should combine both and be refitted as waves come and go; the model fitted on outcomes up to January 2024, which put weight on the journal
rate, ranked fewer later retractions at the very top than the reference count alone.

\paragraph{NIH-funded research.} NIH-funded papers cite already-retracted work as often as other papers in the same
journal and year. About 0.5\% of biomedical papers published 2000--2023 do so: 103,349 papers in the graph, 10,872 of them
NIH-funded. Whether these citations rely on the retracted findings cannot be read from
citation links alone; citation-context resources \citep{joseph2026sanecite,hsiao2021continued} could answer it.

\paragraph{Limitations.} The test is limited to papers in the January 2024 OpenCitations snapshot, which depends on
publishers depositing references with Crossref and is complete only to about September 2023; papers without deposited
references cannot be scored, and 17\% of retracted papers could not be matched to the graph. Retraction Watch is
incomplete: 3,205 papers flagged as retracted in PubMed are missing from the copy used, three quarters of them with
notices from 2024 onwards, so some later retractions are missing from the outcome labels. A retraction is not proof of
misconduct; non-retracted candidates include papers that should have been
retracted, so precision is a lower bound on the share of problematic papers among the flagged ones. The outcome window
is under three years, while the median time from publication to retraction for biomedical papers retracted after January
2024 was about 2.7 years, so
many later retractions are not yet observed. NIH funding is taken from PubMed grant lists, which are incomplete. All
results are associations.
